\hypertarget{appendix-74a0ec818c18}{}
\section*{Linkage and entity ownership}
\section{Linkage and entity ownership}\label{linkage-and-entity-ownership}

\textbf{Curriculum depth:} Introduced · \textbf{First seen:} C007 · \textbf{Latest development:} C007

\subsection{Definition and mechanism}\label{definition-and-mechanism}

Linkage determines when declarations can denote the same entity across scopes or translation units. An unnamed namespace gives internal linkage to suitable names; named modules also introduce module linkage. Linkage is not scope, storage duration or export visibility.

\subsection{Prerequisites}\label{prerequisites}

\begin{itemize}
\tightlist
\item
  \hyperref[appendix-2e4d494c2d94]{Declarations and definitions} --- Declarations must be related to determine whether they denote one entity or distinct private entities.
\end{itemize}

\subsection{Used by}\label{used-by}

Later chunks extend this same node; dependent concepts are listed here as they are integrated.

\subsection{Recall and next depth}\label{recall-and-next-depth}

Give an example, explain why the rule exists, and identify which guarantee would be missing if the prerequisite were ignored.

\textbf{Pending:} Deeper ownership, pointer/reference, generic programming and module rules from later sources. This is written coverage, not demonstrated personal mastery.

\subsection{Source and navigation}\label{source-and-navigation}

\href{Continuous\%20Notes/Computing\%20Foundations\%20-\%20From\%20Bits\%20to\%20C++\%20and\%20Systems\%20APIs\#Many\%20files\%20must\%20describe\%20one\%20coherent\%20program}{C007 teaching chapter} · \hyperref[appendix-fd090d982a15]{Source coverage} · \href{}{Dependency Map (02-Foundations-II, appendix-4ad47bdd269a)}
